% 基礎スライド（日本語）— 「成長がなぜ応力になるか」と「USERMAT とは何か」
%
% ビルド: lualatex slides_kiso_ja.tex （2回）
% 必要パッケージ: luatexja（Debian/Ubuntu なら texlive-lang-japanese、
%   加えて texlive-luatex / texlive-latex-extra）
%
% 使いどころ: 口頭試問の導入、村松研での説明、自分用の確認。
%   slides_muramatsu_ja.tex はこの2点を既知として書いてあるので、
%   相手が初見ならこちらを先に挟む。
%
% CONSTRAINT: 基礎の2点だけ。4ページ。増やさない。
\documentclass[aspectratio=169,11pt]{beamer}
\usepackage{luatexja}
\usetheme{Madrid}
\usecolortheme{seahorse}
\usepackage{amsmath,amssymb}
\usepackage{bm}
\usepackage{xcolor}
\usepackage{tikz}
\usetikzlibrary{arrows.meta,positioning,calc,decorations.pathmorphing}

\definecolor{c3}{RGB}{30,80,160}
\definecolor{cstress}{RGB}{200,50,40}
\setbeamercolor{frametitle}{bg=c3, fg=white}

\title{基礎の2点}
\subtitle{成長がなぜ応力になるのか／USERMAT とは何か}
\author[西岡佳祐]{西岡佳祐}
\date{}

\begin{document}
\maketitle

% --------------------------------------------------------
\begin{frame}{1. 成長がなぜ応力になるのか}
\small 膨らむこと自体は応力を生みません。\textbf{膨らめないこと}が生みます。
\vspace{4pt}
\begin{center}
\begin{tikzpicture}[font=\scriptsize]
  % --- 左：自由に膨らむ ---
  \begin{scope}[xshift=0cm]
    % 膨らんだあと（輪郭だけ）、元のサイズ（破線）はその上に残して見せる
    \draw[fill=c3!8, draw=c3] (-0.3,-0.3) rectangle (1.9,1.9);
    \draw[gray!70, dashed, line width=0.9pt] (0,0) rectangle (1.6,1.6);
    \node[gray!70, font=\tiny, above right] at (0.02,1.62) {元の大きさ};
    % 各辺から外向きに、等方的に広がる
    \foreach \y in {0.4,0.8,1.2} {
      \draw[-{Stealth[length=1.8mm]}, c3] (0,\y) -- ++(-0.28,0);
      \draw[-{Stealth[length=1.8mm]}, c3] (1.6,\y) -- ++(0.28,0);
    }
    \foreach \x in {0.4,0.8,1.2} {
      \draw[-{Stealth[length=1.8mm]}, c3] (\x,0) -- ++(0,-0.28);
      \draw[-{Stealth[length=1.8mm]}, c3] (\x,1.6) -- ++(0,0.28);
    }
    \node[below=1.3cm of {(0.8,0)}, align=center] {自由に膨らむ\\
      \textcolor{c3}{\textbf{応力ゼロ}}};
  \end{scope}
  % --- 右：下がくっついている ---
  \begin{scope}[xshift=6.4cm]
    \draw[fill=cstress!12, draw=cstress] (-0.3,0) rectangle (1.9,1.9);
    % 固着面
    \draw[line width=1.2pt, black] (-0.5,0) -- (2.1,0);
    \foreach \x in {-0.4,-0.1,...,2.0}
      \draw[black] (\x,0) -- ++(-0.18,-0.18);
    \node[below=2mm of {(0.8,-0.15)}, font=\scriptsize] {歯の表面（動かない）};
    % 押し合いの矢印
    \foreach \y in {0.45,0.95,1.45} {
      \draw[-{Stealth[length=2mm]}, cstress, thick] (-0.05,\y) -- ++(0.45,0);
      \draw[-{Stealth[length=2mm]}, cstress, thick] (1.65,\y) -- ++(-0.45,0);
    }
    \node[below=1.3cm of {(0.8,0)}, align=center] {くっついたまま膨らむ\\
      \textcolor{cstress}{\textbf{内部に応力}}};
  \end{scope}
\end{tikzpicture}
\end{center}
\vspace{10pt}
{\footnotesize パン生地が型の中で膨らむのと同じです。型より大きくなろうとすると、
型が押し返す。バイオフィルムは歯の表面に固着したまま膨らむので、
\textbf{自由に膨らめない分が応力になります}。}
\end{frame}

% --------------------------------------------------------
\begin{frame}{$\alpha$ と $\bm{F}_g=(1+\alpha)\bm{I}$ の意味}
\small 式は1つだけ覚えれば足ります。
\vspace{6pt}
\begin{center}
\colorbox{c3!10}{\parbox{0.8\linewidth}{\centering
\large $\bm{F}_g=(1+\alpha)\,\bm{I}$\\[4pt]
\small 「全方向に $(1+\alpha)$ 倍に\textbf{膨らみたい}」}}
\end{center}
\vspace{8pt}
\begin{itemize}\small
  \item $\alpha$ は\textbf{膨らみたい量}。$\alpha=0.05$ なら「5\% 大きくなりたい」
  \item $\alpha$ が大きいほど押し合いが強く、応力も大きい
  \item $\bm{I}$ は単位行列。\textbf{どの方向にも同じだけ}膨らむ、という意味
  \item $\alpha$ は細菌の組成から決まる
        $\Rightarrow$ \textbf{菌の種類が変われば応力が変わる}
\end{itemize}
\vspace{6pt}
{\footnotesize 「膨らみたい」と「実際に膨らんだ」は別物です。
実際の変形は、まわりが許す分しか起きません。
その\textbf{差}が応力になります。}
\end{frame}

% --------------------------------------------------------
\begin{frame}{2. USERMAT とは何をするものか}
\small ANSYS は計算中、各点・各ステップで\textbf{何万回も}同じことを聞いてきます。
\vspace{4pt}
\begin{center}
\begin{tikzpicture}[font=\small, node distance=8mm,
  b/.style={draw, rounded corners=3pt, minimum height=11mm, text width=30mm,
            align=center},
  ar/.style={-{Stealth[length=2.4mm]}, thick}]
  \node[b, fill=gray!10] (ans) {ANSYS\\（構造解析）};
  \node[b, fill=c3!18, draw=c3, very thick, right=36mm of ans] (um)
       {\textbf{USERMAT}\\自分で書いた材料則};
  \draw[ar, gray] ([yshift=3mm]ans.east) -- node[above, font=\scriptsize, align=center]
       {この点はこれだけ\\変形している} ([yshift=3mm]um.west);
  \draw[ar, c3] ([yshift=-3mm]um.west) -- node[below, font=\scriptsize, align=center]
       {応力はこれです} ([yshift=-3mm]ans.east);
\end{tikzpicture}
\end{center}
\vspace{2pt}
\begin{itemize}\small
  \item 普通のゴムや金属なら、ANSYS が答えを最初から持っている
  \item しかし「\textbf{膨らもうとするゴム}」は持っていない
  \item $\Rightarrow$ その答えを返す部品を自分で書いた。\textbf{これが修論の本体}
\end{itemize}
\vspace{4pt}
{\footnotesize ANSYS 本体には一切手を入れていません。
「聞かれたら答える」部品を差し替えているだけです。
だから ANSYS のバージョンにも依存しません。}
\end{frame}

\end{document}
